\problem{}

This problem is mainly another application of Noether's theorem. But, it serves as a nice example of a symmetry transformation 
which changes the space-time coordinates with an non unit determinant. This scale symmetry is a particular case of the more 
generic conformal transformations which is a whole area of research in QFT, it's also a nice exemple of an anomalous symmetry. 
Despite it being a symmetry of the classical theory, in most of the cases the scale symmetry is absent in the quantum theory, 
we probably will see this by the end of the course when we introduce beta functions and the anomalous mass dimension. There are a 
few known interacting theories which indeed retains scale/conformal symmetry after quantization. 

Again we will use \cref{eq:conservedcurrentnoether}, in which we just need to know the transformation law,
\begin{align}
    \begin{cases}
        x&\rightarrow x' = \e^a x\\
        \phi\qty(x)&\rightarrow \phi'\qty(x')=\e^{-d_\phi a}\phi\qty(x)
    \end{cases}\Rightarrow \dv{{x'}^\mu}{a}\eval_{a=0} = x^\mu,\ \ \ \dv{\phi'\qty(x)}{a}\eval_{a=0}=-d_\phi\phi\qty(x) -\partial_\mu\phi\qty(x)x^\mu
\end{align}
it's interesting to notice that the scale transformation actually ``measures'' the mass dimension of the objects being transformed, $x$ has 
mass dimension $\qty[x]=-1$, by consistency our field has mass dimension $\qty[\phi]=1=d_\phi$ which fixes the transformation law as $\phi'(x')\rightarrow \e^{-a\qty[\phi]}\phi\qty(x)$. 
Also, is nice to check that $\qty[\lambda]=0$. We then start by showing that it's a symmetry,
\begin{align}
    S'_a\qty[\phi']&= \int\dd[4]{x'}\qty[-\frac12\partial'_\mu\phi'\qty(x')\partial'^\mu\phi'\qty(x')-\frac{\lambda}{4!}\phi'^4\qty(x')]\\
    S'_a\qty[\phi']&= \int\dd[4]{x}\e^{4a}\qty[-\frac12\e^{-2a-2d_\phi a}\partial_\mu\phi\qty(x)\partial^\mu\phi\qty(x)-\e^{-4d_\phi a}\frac{\lambda}{4!}\phi^4\qty(x)]\\
    S'_a\qty[\phi']&= \int\dd[4]{x}\qty[-\frac12\e^{2a-2d_\phi a}\partial_\mu\phi\qty(x)\partial^\mu\phi\qty(x)-\e^{4a-4d_\phi a}\frac{\lambda}{4!}\phi^4\qty(x)]
\end{align}
setting $d_\phi=1$,
\begin{align}
    S'_a\qty[\phi']&= \int\dd[4]{x}\qty[-\frac12\partial_\mu\phi\qty(x)\partial^\mu\phi\qty(x)-\frac{\lambda}{4!}\phi^4\qty(x)]=S[\phi]
\end{align}
what settles down as a symmetry.

Now, let's compute the conserved charge, again, we use the above derived infinitesimal transformations in \cref{eq:conservedcurrentnoether},
\begin{align}
    D^\mu &=\pdv{\mathcal L}{\partial_\mu\phi\qty(x)}\dv{\phi'\qty(x)}{a}\eval_{a=0}+\mathcal L\dv{{x'}^\mu}{a}\eval_{a=0}\\
    D^\mu &=-\pdv{\mathcal L}{\partial_\mu\phi\qty(x)}\qty(d_\phi\phi\qty(x) +\partial_\alpha\phi\qty(x)x^\alpha)+\mathcal Lx^\mu\\
    D^\mu &=-d_\phi\pdv{\mathcal L}{\partial_\mu\phi\qty(x)}\phi\qty(x)-\pdv{\mathcal L}{\partial_\mu\phi\qty(x)}\partial_\alpha\phi\qty(x)x^\alpha+\mathcal Lx^\mu\\
    D^\mu &=-d_\phi\pdv{\mathcal L}{\partial_\mu\phi\qty(x)}\phi\qty(x)+x^\alpha\qty(-\pdv{\mathcal L}{\partial_\mu\phi\qty(x)}\partial_\alpha\phi\qty(x)+\mathcal L\tensor{g}{^\mu_\alpha})\\
    D^\mu &=\partial^\mu\phi\phi+x^\alpha\tensor{T}{^\mu_\alpha}\label{eq:dilatoncurrentunimproved}
\end{align}
in the last line we used $d_\phi=1$ and identified the quantity multiplying $x^\alpha$ as the energy-momentum tensor. This format is more insightful, but for explicit 
calculations is better to expand into the explicit $\phi$ dependence,
\begin{align}
    D^\mu &=\partial^\mu \phi\qty(d_\phi\phi +\partial_\alpha\phi x^\alpha)+\qty(-\frac12\partial^\alpha\phi\partial_\alpha\phi-\frac{\lambda}{4!}\phi^4)x^\mu
\end{align}

Again, as a sanity check we show it is divergenceless,
\begin{align}
    \partial_\mu D^\mu &=\partial_\mu \partial^\mu \phi\qty(d_\phi\phi +\partial_\alpha\phi x^\alpha)+\partial^\mu \phi\qty(d_\phi\partial_\mu \phi +\partial_\mu\partial_\alpha\phi x^\alpha+\partial_\alpha\phi\partial_\mu x^\alpha)\nonumber\\
    &\quad\quad\quad+\qty(-\partial^\alpha\phi\partial_\mu\partial_\alpha\phi-\frac{\lambda}{3!}\phi^3\partial_\mu\phi)x^\mu+\qty(-\frac12\partial^\alpha\phi\partial_\alpha\phi-\frac{\lambda}{4!}\phi^4)\partial_\mu x^\mu\\
    \partial_\mu D^\mu &=\Box\phi\qty(d_\phi\phi +\partial_\alpha\phi x^\alpha)+\partial^\mu \phi\qty(d_\phi\partial_\mu \phi +\partial_\mu\partial_\alpha\phi x^\alpha+\partial_\alpha\phi\tensor{g}{_\mu^\alpha})\nonumber\\
    &\quad\quad\quad+\qty(-\partial^\alpha\phi\partial_\mu\partial_\alpha\phi-\frac{\lambda}{3!}\phi^3\partial_\mu\phi)x^\mu+4\qty(-\frac12\partial^\alpha\phi\partial_\alpha\phi-\frac{\lambda}{4!}\phi^4)\\
    \partial_\mu D^\mu &=\qty(\Box\phi-\frac{\lambda}{3!}\phi^3)x^\alpha\partial_\alpha\phi+\qty(d_\phi\Box\phi-\frac{\lambda}{3!}\phi^3)\phi+\partial^\mu \phi\qty(d_\phi\partial_\mu \phi +\partial_\mu\partial_\alpha\phi x^\alpha+\partial_\mu\phi)\nonumber\\
    &\quad\quad\quad-\partial^\alpha\phi\partial_\mu\partial_\alpha\phi x^\mu-2\partial^\alpha\phi\partial_\alpha\phi\\
    \partial_\mu D^\mu &=\qty(\Box\phi-\frac{\lambda}{3!}\phi^3)x^\alpha\partial_\alpha\phi+\qty(\Box\phi-\frac{\lambda}{3!}\phi^3)\phi+\qty(d_\phi-1)\phi\Box\phi+\qty(d_\phi-1)\partial^\mu \phi\partial_\mu \phi
\end{align}
using the equation of motion, \begin{align}\Box\phi=\frac{\lambda}{3!}\phi^3\end{align}
we get,
\begin{align}
    \partial_\mu D^\mu &=\qty(d_\phi-1)\phi\Box\phi+\qty(d_\phi-1)\partial^\mu \phi\partial_\mu \phi
\end{align}
finally setting $d_\phi=1$, 
\begin{align}
    \partial_\mu D^\mu &=0
\end{align}

Just as a curiosity, \cref{eq:dilatoncurrentunimproved} can be rewritten as,
\begin{align}
    D^\mu = \partial^\mu\qty(\frac12 \phi^2)+x^\alpha \tensor{T}{^\mu_\alpha}
\end{align}
the first term is a total derivative, and usually total derivative contributions to the current can be ignored. 
For general theories the scale current has the form
\begin{align}
    D^\mu \sim x^\alpha \tensor{\Theta}{^\mu_\alpha} +\textnormal{ total derivative}
\end{align}
thus, it's conservation is dependent of, assuming translational invariance,
\begin{align}
    \partial_\mu D^\mu \sim \partial _\mu x^\alpha \tensor{\Theta}{^\mu_\alpha}+x^\alpha \partial_\mu\tensor{\Theta}{^\mu_\alpha}\\
    \partial_\mu D^\mu \sim  \tensor{\Theta}{^\mu_\mu}\stackrel{?}{=}0
\end{align}
as long as the trace of the energy-momentum tensor is zero, up to total derivatives, the theory under consideration is 
scale invariant. This is all classical, under quantization the matter is way more intricate. In four dimensions and under the usual assumptions 
of unitarity, Poincaré invariance, locality, causality, all known interacting scale invariant theories are conformal ones. Example of 
non-interacting scale invariant theories are easy, any massless theory, in fact for pure electromagnetism we have $\tensor{\Theta}{^\mu_\mu}=0$. 
For interacting theories, QCD, or any Yang-Mills are classical scale invariant, but, under quantization aren't. This means the 
quantum energy momentum tensor has non-zero trace, despite the classical one being zero. Due to dimensional reasons the non-zero 
contribution generated by quantization comes accompanied by positive powers of $\hbar$,
\begin{align}
    \tensor{{\hat {\Theta}}}{^\mu_\mu}=\hbar^n C
\end{align}
this phenomenon is usually called \textit{Trace anomaly}. It plays an important role in String Theory also. Just to mention, 
one of the known interacting quantum conformally invariant theories is a so called $\mathcal N=4$ Super-Yang-Mills.